\documentclass[UTF8,11pt]{ctexart}
\usepackage{amsmath}
\usepackage{amssymb}
\usepackage[a4paper,margin=2.2cm]{geometry}
\usepackage{booktabs}
\usepackage{array}
\usepackage{longtable}
\usepackage{enumitem}
\usepackage{xcolor}
\usepackage{url}
\usepackage{array}
\usepackage[hidelinks]{hyperref}
\hypersetup{colorlinks=true,linkcolor=blue,urlcolor=blue,citecolor=blue}
\newcommand{\codeinl}[1]{\texttt{#1}}

\title{\bfseries 全国量化基金人才需求与供给调查\\[6pt]\large 数据均来自官方与可查证的公开渠道}
\author{量化投研调查笔记}
\date{}

\begin{document}

\maketitle

\begin{abstract}
本报告在严格遵循“数据可溯源、口径可核验”的前提下，梳理国内量化基金行业的人才需求与高校相关专业的人才供给。与一般行业报道不同，本报告\textbf{不采信无法溯源至官方或权威数据库的数字}：凡是只能由第三方商业机构估算或仅见于行业随口数据者，或删除、或降级标注、或明确声明“无官方统计”。需求侧以中国证券投资基金业协会（AMAC）官方口径为准，供给侧以教育部《普通高等学校本科专业目录》与教育部阳光高考专业库为准。报告的核心发现是：\textbf{官方层面不存在“量化私募从业人员”这一细分统计口径}，任何“量化从业者 X 万人”的说法都是推算或估算，而非权威数据。据此，本报告对“金融工程 / 金融数学是否值得攻读”仅给出基于可核验证据的条件性判断，并明确标注其推理依赖第三方数据之处。
\end{abstract}

\tableofcontents
\newpage

\section{引言：这个问题值得认真查一查}

\subsection{为什么要关心}
一个看似简单的问题——“全国量化基金有多少从业者、高校相关专业能供给多少人”——常常被各种行业报道随口答成一个很大的数字（“几万人”“几十万人”）。\textbf{但这些数字大多找不到官方出处。}这份调查的价值，就是把“能查证的事实”和“查不到的估算”严格分开，让读者看清：哪些数字有权威出处、哪些只是行业的估算。

\subsection{我们是怎么查的}
做这份调查，我们坚持三条原则：
\begin{enumerate}[itemsep=2pt]
\item \textbf{数字要有出处}：把数据分成“官方发布”“官方平台或权威机构”“第三方调研或商业估算”三档，用的时候逐一标清楚是哪一档。
\item \textbf{查不到出处的不用}：凡是查不到可靠出处的数字，不写进正文结论，只在后面说明为什么不能用。
\item \textbf{留痕可复核}：每个数字都写明出处机构、文件名称、查询链接和查询日期，方便读者自己再去核对。
\end{enumerate}

\section{数据来源与口径}

\begin{table}[ht]
\centering
\small
\begin{tabular}{@{}p{3.0cm}p{4.4cm}p{4.3cm}@{}}
\toprule
\textbf{层级} & \textbf{性质} & \textbf{本报告的使用方式} \\
\midrule
A 官方一手 & 监管机构/政府部门统计或明文条款 & 作为正文数据的主要依据 \\
B 官方平台/权威机构 & 学信网·阳光高考、高校官方、行业自律组织 & 作为补充依据，注明出处 \\
C 第三方调研/商业估算 & 商业公司白皮书、媒体转述 & 一般\textbf{不进入正文}，仅在标注后参考 \\
\bottomrule
\end{tabular}
\caption{数据来源分层说明}
\end{table}

\noindent 文中数据的查询日期统一为 \textbf{2026-09-22}。

\section{需求侧：量化从业者到底有多少}

\subsection{先说结论：官方没有“量化从业者”这个统计}

\textbf{核心发现}：中国证券投资基金业协会（AMAC）的从业人员管理平台\textbf{并未按“量化策略/量化私募”拆分公布从业人员数据}。协会统计的是“私募基金管理人”这一整类。换句话说，\textbf{“全国量化从业者有多少”在官方层面找不到答案}，任何具体数字都来自第三方推算。\footnote{经查证，AMAC 年报的“从业人员”章节仅按私募基金整类披露，未拆分量化细分。}

一个旁证是：业界测算“量化私募从业人数”时，仍需拿“私募证券基金从业人员总数 $\times$ 一个量化的占比”去估算，而不是直接查一张官方表——这反过来说明官方没有量化细分口径。\footnote{见第三方《2021 年中国量化投资白皮书》以“中基协私募证券投资基金从业人员 $\times$ 量化占比”推算之做法。}

\subsection{能查到的第一层：私募管理人的官方从业数据}

目前能从 AMAC 官方查到的、较新的口径是 \textbf{2022 年度}数据（AMAC《2022 年私募基金登记备案综述》，经济参考报转引 AMAC 口径）：
\begin{itemize}[itemsep=2pt]
\item 私募基金管理人在\textbf{从业人员管理平台完成注册的全职员工}为 \textbf{18.30 万人}；
\item 其中\textbf{取得基金从业资格}的员工为 \textbf{15.69 万人}。\footnote{中国证券投资基金业协会《2022 年私募基金登记备案综述》，经济参考报转引，2023-01。链接：\url{http://www.jjckb.cn/2023-01/03/c_1310688104.htm}。另见 AMAC《中国证券投资基金业年报（2022）》“私募投资基金—从业人员整体情况”章节（第 260 页左右）：\url{https://www.amac.org.cn/sjtj/tjbg/nb/202312/P020231204601450460524.pdf}。查询日期：2026-09-22。}
\end{itemize}

需要特别说明：这 18.30 万人是\textbf{全部私募基金}（含股权/创投/证券等）的从业人员，\textbf{不区分量化与否}，也\textbf{不等于“量化从业者”}。更新的官方月报/综述以 AMAC 官网“统计数据”栏目实时发布为准。

\subsection{能查到的第二层：监管的底线要求}

“私募基金管理人需要有一定数量的正式员工，且员工要缴纳社保、签劳动合同”——这是\textbf{写在监管规则里的明文要求}，来源清楚：
\begin{itemize}[itemsep=2pt]
\item \textbf{条款}：《私募投资基金登记备案办法》（AMAC，2023 年 5 月施行）第八条第一款第（四）项要求管理人“专职员工不少于 5 人”；配套《私募基金管理人登记指引第 1 号——基本经营要求》第九条明确：\textbf{“专职员工”是指与管理人签订劳动合同并缴纳社保的正式员工}。\footnote{原文：《私募投资基金登记备案办法》PDF：\url{https://www.amac.org.cn/fwdt/wyb/jgdjhcpbeian/smjjglrdjhcpba/fwzn/202206/P020231129532332167341.pdf}；《私募基金管理人登记指引第 1 号》全文：\url{https://fg.amac.org.cn/governmentrules_3854/zcgz_zlgz/zlgz_smjj/zlgz_smjj_glrdj/202303/t20230327_24588.html}；登记申请材料清单（含社保缴费记录要求）：\url{https://www.amac.org.cn/fwdt/wyb/jgdjhcpbeian/smjjglrdjhcpba/fwzn/202309/P020231126371388739564.pdf}。查询日期：2026-09-22。}
\end{itemize}

这条规则的意义在于：监管以“${}\ge 5$ 名专职 + 社保 + 劳动合同”作为一家照牌机构的合规底线，可以用它来理解机构从业规模的下限，但它\textbf{不等于行业实际从业人数}。

\subsection{头部（百亿）私募的规模：家数是官方的，从业数是第三方}

很多人关注的“百亿量化私募有多少人”，需要把“家数”和“从业人数”分开看，两者权威层级完全不同：

\begin{itemize}[itemsep=2pt]
\item \textbf{家数是官方口径}：证监会主席吴清 2026-06 在中基协致辞中披露，私募行业已有“\textbf{十余家千亿级、400 余家百亿级管理人}”（为累计口径，含股权/创投等全类型，未按量化拆分）。\footnote{证监会主席吴清在中国证券投资基金业协会第四届会员代表大会上的致辞，发布于 AMAC 官网“证监要闻”：\url{https://www.amac.org.cn/xwfb/zjyw/202606/t20260606_27783.html}。查询日期：2026-09-22。}
\item \textbf{从业人数的官方口径却查不到}：AMAC/证监会\textbf{只公布“百亿私募”的家数，从不公布这些头部机构的从业人数汇总}。官方月报仅按私募证券（7,247 家）/股权创投（11,103 家，2026-07 末）分类，无头部从业人数；中基协从业人员公示平台虽能逐家查询，但\textbf{没有官方汇总数}。
\end{itemize}

因此“百亿量化私募有多少从业者”，公开能查到的\textbf{只有第三方统计}。例如私募排排网口径的“\textbf{74 家百亿量化私募，员工约 4,495 人、平均每家约 61 人}”（2026 年 7 月时点）：\footnote{数据来源：私募排排网统计，经新浪财经等转载《百亿量化私募持续高薪抢人才！年薪80万起步！还有股权等福利》：\url{https://cj.sina.com.cn/articles/view/5996682579/1656e1d53001018k20}。查询日期：2026-09-22。该数字为\textbf{第三方统计}，非官方发布；且该时点口径与同期“71 家”“85 家”等版本并存、并不一致。}

\textbf{处置说明}：这份调查以“数据可溯源”为原则，原则上不采信无官方出处的数字。但“头部私募从业规模”是读者高度关注的信息，且没有更好的官方替代，因此本调查\textbf{破例采用该第三方口径，但明确标注它是私募排排网的统计、时点截至 2026-07，并注明口径存在多版本差异}。读者引用时应注明其为第三方数据、而非官方。

\textbf{更可信的替代}：若需要更可信的头部从业数，唯一相对可靠的做法是去\textbf{中基协从业人员公示平台}逐家自行汇总（可核验，但非官方汇总数）。

\subsection{依然不能采信的“量化从业规模”估算}

下面这些说法\textbf{查不到任何统计出处}，这份调查不采用：
\begin{itemize}[itemsep=2pt]
\item “机构量化团队 1--2 万人”“个人量化交易 10--20 万人”——都是\textbf{行业经验性的估算}，没有统计出处；
\item “合规缴纳社保的量化从业者约 3,000--4,500 人”——这是此前的推算值，\textbf{本调查予以撤回}，因为它建立在一个量化的占比假设上，AMAC 没有这样的官方依据。
\end{itemize}

\subsection{需求侧小结}

把能查到的官方信息放一起，需求侧能下的\textbf{可靠结论}是：
\begin{quote}
官方不公布“量化从业人员”的细分数据；目前能查证的可靠官方信息包括：私募基金整类从业人数（2022 年末全职注册 18.30 万人、具从业资格 15.69 万人）、“专职员工 $\ge 5$ 名且缴纳社保”的监管底线，以及“百亿私募累计约 400 余家”的官方家数。头部（百亿）私募的\textbf{从业人数}官方不公布，公开仅有私募排排网的第三方统计（如 2026-07 按 74 家、约 4,495 人，多口径并存）。所以\textbf{“量化从业者的精确规模”无法从官方数据直接得出}。
\end{quote}

\section{供给侧：高校相关本科专业}

\subsection{专业目录：这些专业是官方承认的}

教育部《普通高等学校本科专业目录》是供给侧最权威的一手来源。现行版本（2025 年）收录以下相关专业：\footnote{教育部《普通高等学校本科专业目录（2025 年）》PDF：\url{http://www.moe.gov.cn/srcsite/a08/moe_1034/s4930/202504/W020250422312780837078.pdf}；发文号教高函〔2025〕3 号。查询日期：2026-09-22。}
\begin{itemize}[itemsep=2pt]
\item \textbf{金融工程}（020302，基本专业）；
\item \textbf{金融数学}（020305T，特设专业）；
\item \textbf{精算学}（020308T，可授经济学或理学）；
\item \textbf{互联网金融}（020309T，特设专业，2016 年增设）；
\item \textbf{金融科技}（020310T，特设专业，2017 年增设）；
\item \textbf{金融审计}（020311TK，特设+国家控制布点专业）。
\end{itemize}
（增设年份以教育部《普通高等学校本科专业目录（2020 年版）》为准。）

\textbf{结论}：教育部目录说明——量化相关培养并不是一个独立的“量化专业”，而是\textbf{由金融工程 / 金融数学 / 金融科技 / 精算学等承接}。这也印证了多数人的认知：它属于“金融学大类下的数量方向”。

\subsection{毕业生规模：有官方口径}

教育部阳光高考 / 学信网专业库提供了\textbf{官方}的毕业生规模区间（统计截止 2025-12-31）：
\begin{itemize}[itemsep=2pt]
\item \textbf{金融工程}：全国普通高校毕业生规模约 \textbf{22,000--24,000 人}；\footnote{阳光高考专业库·金融工程：\url{https://gaokao.chsi.com.cn/zyk/zybk/wap/detail/73381107}。查询日期：2026-09-22。}
\item \textbf{金融科技}：全国普通高校毕业生规模约 \textbf{3,000--3,500 人}；\footnote{阳光高考专业库·金融科技：\url{https://gaokao.chsi.com.cn/zyk/zybk/wap/detail/2jgwq0d4turhi4og}。查询日期：2026-09-22。}
\item 更宽的\textbf{金融学专业}本科毕业生规模约 \textbf{60,000--65,000 人}（另有口径称 70,000--75,000，版本区间需注明）。\footnote{阳光高考专业库·金融学：\url{https://gaokao.chsi.com.cn/zyk/zybk/viewPage/vgloy4xrpyaztqbf}。两处数字存在口径年份差异，引用时需要注明。}
\end{itemize}
需注意：官方只给\textbf{毕业生规模区间}，没有逐年精确数，也没有分专业的就业落实率。

\subsection{开设高校数量：没有官方统一数字}

\textbf{“全国有多少所高校开设金融工程 / 金融数学 / 金融科技”——官方没有给出一个统一的汇总数字。}教育部并未按专业公开开设院校总数，只有阳光高考专业库可以逐校查询，其余多是第三方商业统计：
\begin{itemize}[itemsep=2pt]
\item 阳光高考（学信网）专业解读称“开设金融工程的院校有 260 多所”；\footnote{阳光高考专业解读·金融工程：\url{https://gaokao.chsi.com.cn/zyk/zybk/zyjd/viewPage/tcl1qp6hpacg3a9m}。查询日期：2026-09-22。}
\item 金融科技开设院校数，各家第三方口径不一（如清华五道口报告“110 余所”\footnote{清华大学五道口金融学院《金融科技人才供需调研报告（2024）》：\url{https://thuifr.pbcsf.tsinghua.edu.cn/}（报告 PDF 内“截至 2024 年 6 月开设金融科技本科院校 110 余所”）。查询日期：2026-09-22}、掌上高考“149 所”\footnote{掌上高考（第三方）：\url{https://www.gaokao.cn/gk-mb/31/8322-0}。查询日期：2026-09-22}等）。
\end{itemize}
\textbf{小结}：这类数字都是第三方统计，各家口径不同，\textbf{不能当作官方数据引用}。

\subsection{代表高校：能查到培养方式，查不到具体人数}

以中央财经大学为例，学校官网和培养方案能查到\textbf{培养机制}（属官方一手），但\textbf{没有公开每年招生 / 升学的具体人数}：
\begin{itemize}[itemsep=2pt]
\item 金融工程 2002 年设立，是全国\textbf{首批五个金融工程本科专业点之一}；与金融科技同属一个招生大类，\textbf{大二专业分流}。\footnote{中央财经大学金融学院·金融工程专业页：\url{https://sf.cufe.edu.cn/jxxm/bk/jrgczy.htm}；教务处培养方案 PDF：\url{https://jwc.cufe.edu.cn/PDP2024/13_jinronggongcheng.pdf}。查询日期：2026-09-22。}
\item 金融科技 2019 年首届本科招生。\footnote{中央财经大学金融学院学科简介：\url{https://sf.cufe.edu.cn/xygk/xyjj/xkjj.htm}。查询日期：2026-09-22。}
\item \textbf{不足}：该校官方\textbf{没有公开}每年的招生总额、分流规模上限和升学深造率，因此本调查不采信“央财金融工程年招 19--22 人”“深造率逾七成”这类没有官方文件佐证的数字。
\end{itemize}

\subsection{供给侧小结}

\begin{quote}\sloppy
供给侧能查证的官方事实：教育部目录有金融工程 / 金融数学 / 金融科技等专业（020302、020305T、020310T 等）；阳光高考给官方口径：金融工程年毕业 2.2--2.4 万、金融科技 0.3--0.35 万。查不到官方出处的（我们不用）：开设高校总数、各校招生人数、升学深造率的具体数值。
\end{quote}

\section{供需匹配：能查证的事实 vs 查不到的估算}

\subsection{供需两侧的数据：存量与增量，不可直接比}

把需求侧与供给侧能从官方查证的数据放到一起，必须先认清一个关键差异：\textbf{需求侧是“存量”，供给侧是“流量”}。
\begin{table}[ht]
\centering
\small
\begin{tabular}{@{}>{\raggedright\arraybackslash}p{3.9cm}>{\raggedright\arraybackslash}p{2.6cm}>{\raggedright\arraybackslash}p{1.6cm}>{\raggedright\arraybackslash}p{3.6cm}@{}}
\toprule
\textbf{维度（官方口径）} & \textbf{数字} & \textbf{口径性质} & \textbf{出处} \\
\midrule
私募基金管理人全职注册从业（2022 末） & 18.30 万人 & 存量（在岗总人数） & AMAC 官方综述 \\
其中具从业资格（2022 末） & 15.69 万人 & 存量（在岗总人数） & AMAC 官方综述 \\
百亿量化私募从业（2026-07） & 约 4,495 人 & 存量（头部在岗）；\textbf{第三方统计} & 私募排排网 \\
金融工程本科年毕业生 & 2.2--2.4 万 & \textbf{年增量（当年毕业）} & 教育部阳光高考 \\
金融科技本科年毕业生 & 0.3--0.35 万 & \textbf{年增量（当年毕业）} & 教育部阳光高考 \\
监管：专职员工 $\ge 5$ 名且缴社保 & 条款 & 底线要求 & AMAC 自律规则 \\
\bottomrule
\end{tabular}
\caption{供需两侧官方数据：存量与增量口径对照}
\end{table}

\textbf{为什么这些存量数字不能与供给流量直接比较}：需求侧的量，无论是私募基金全体在岗的 18.30 万人，还是头部（百亿）量化私募的约 4,495 人，都是\textbf{截至某时点的存量人数}；而金融工程 + 金融科技的本科年毕业生（约 2.5--2.8 万），是\textbf{当年新进入劳动力市场的流量}。拿“存量的总盘子”去比“一年的新增”，既不能说“供给填满了需求”，也不能说“供给不足”——因为要判断供需，缺的是\textbf{“行业每年新增岗位（流入）”这个流量口径的需求数据}，而这恰恰是官方没有公布的。（需再次强调：头部从业数 4,495 人为\textbf{第三方统计}，仅作存量量级参考。）

\textbf{把两侧数据放到同一条时间轴上，另一层不可比更直观}：需求侧只有“某时点上有多少人”（存量），供给侧只有“某年份毕业了多少人”（流量），两侧的时间点交错出现、口径各不相同，永远进不了同一个“供需比”的分母。
\begin{table}[htbp]
\centering
\small
\renewcommand{\arraystretch}{1.3}
\begin{tabular}{@{}>{\raggedright\arraybackslash}p{1.9cm}>{\raggedright\arraybackslash}p{6.4cm}>{\raggedright\arraybackslash}p{6.2cm}@{}}
\toprule
\textbf{时间} & \textbf{需求侧（存量 / 时点口径）} & \textbf{供给侧（流量 / 年度口径）} \\
\midrule
2002 & — & 金融工程专业设立（全国首批五个本科专业点之一，央财） \\
2016 & — & 互联网金融专业增设（教育部目录） \\
2017 & — & 金融科技专业增设（教育部目录） \\
2019 & — & 金融科技首届本科招生（央财） \\
2021 & 券商从业者硕士占比 26.81\%（中证协研究） & — \\
2022 末 & 私募基金管理人全职注册从业 18.30 万；其中具从业资格 15.69 万（AMAC 官方综述） & — \\
2024 末 & 券商从业者硕士占比升至 34.06\%（中证协研究） & — \\
2025-12-31 截止 & — & 官方毕业生规模区间：金融工程 2.2--2.4 万、金融科技 0.3--0.35 万（阳光高考） \\
2026-07 & 百亿量化私募从业约 4,495 人（\textbf{第三方}：私募排排网） & — \\
\bottomrule
\end{tabular}
\caption{供需两侧数据时间轴对照：需求侧是时点存量，供给侧是年度流量}
\end{table}

这张时间轴的读法是：两侧各自有可查证的事实，但\textbf{没有任何一个年份同时给出需求侧的新增岗位流量与供给侧的毕业生流量}——决定供需比的分子（新增需求）从未与分母（新增供给）同框出现，因此任何“供需比”数字都没有官方依据，只有口径对齐的“当年新增岗位”数据出现后，这个问题才可能被回答。

\textbf{能下的可靠判断}：
\begin{enumerate}[itemsep=2pt]
\item 需求侧存量（无论是 18.30 万全体在岗，还是头部约 4,495 人）与供给侧流量（每年 2.5--2.8 万毕业生）\textbf{不属于同一口径}，本调查\textbf{不做任何“供大于求 / 供不应求”的数量结论}，因为缺少官方的“年度新增岗位”流量数据作为可比基准。
\item 更关键的是：\textbf{“量化”这个标签在官方两边都不存在}（需求侧没有量化的细分，供给侧没有量化的专业），所以“量化人才供需是否匹配”\textbf{无法用官方数据直接回答}，只能落到“金融工程 / 金融数学毕业生进入金融行业”这个间接问题上。
\end{enumerate}

\subsection{查不到、证不了的说法}

任何“量化岗位缺口 X 万人”“量化从业者 XX 万人”“供需比 X”这样的数字，这份调查\textbf{都不作背书}，除非给出能查证的官方或权威统计依据。也就是说，\textbf{行业里常见的“整体供大于求 / 供不应求”结论，我们既无法证实也无法证伪}。

\section{回到问题：金融工程 / 金融数学到底值不值得读}

在“查不到出处的数字不用”的约束下，这份调查\textbf{不给}具体的录取率、薪资、就业率结论（因为官方没有这些数据）。我们只基于能查证的事实，给一个\textbf{分情况的判断框架}：

\subsection{几条能查证的客观前提}
\begin{itemize}[itemsep=2pt]
\item \textbf{专业是官方承认的，且培养规模不小}：教育部目录列有金融工程、金融数学等专业，金融工程年毕业生 2.2--2.4 万（官方口径），说明这是\textbf{大规模培养}的方向，而不是冷门实验班。
\item \textbf{金融行业学历门槛确实在上升，这有官方旁证}：中国证券业协会的研究显示，截至 2024 年末券商从业者硕士及以上占比 \textbf{34.06\%}（2021 年是 26.81\%）\footnote{中国证券业协会研究（经证券时报披露）：\url{https://www.stcn.com/article/detail/3149883.html}。查询日期：2026-09-22。}——需要说明的是，这统计的是\textbf{券商整体}（含非量化岗位），属行业自律组织发布。金融类专业如果本科直接就业，客观上会面对学历竞争上升的局面。
\item \textbf{量化岗位的学历结构只有第三方数据}：第三方（朝阳永续/华泰）称量化从业者中博士 11\%、硕士 65\%——这是\textbf{商业机构调研}，本调查仅作参考，不作为官方事实。\footnote{来源：证券时报引第三方《2021 中国量化私募白皮书》关于学历结构：\url{https://www.stcn.com/stock/djjd/202209/t20220906_4843317.html}、\url{https://file.joinquant.com/a265153cab9feb4c03bc9305abe9de37/2021\%E5\%B9\%B4\%E5\%BA\%A6\%E4\%B8\%AD\%E5\%9B\%BD\%E9\%87\%8F\%E5\%8C\%96\%E6\%8A\%95\%E8\%B5\%84\%E7\%99\%BD\%E7\%9A\%AE\%E4\%B9\%A6.pdf}。查询日期：2026-09-22。}
\end{itemize}

\subsection{分人群的判断}
\begin{enumerate}[itemsep=2pt]
\item \textbf{如果目标是想进证券 / 基金 / 金融机构}：金融工程这类专业是\textbf{官方承认、规模培养}的对口方向，但学历硬门槛客观存在（券商硕士占比 34\% 且在上升），所以“本科毕业就靠这个专业直接就业”\textbf{现实上难度较高}；继续读研或读博是更匹配的路。
\item \textbf{如果目标只是泛金融或非投研岗位}：数理背景是加分项，专业选择的影响有限。
\item \textbf{如果对“量化核心岗”抱有具体期待}：由于薪资、录取率都没有官方数据，本调查\textbf{无法给出明确结论}，只能提醒：这类判断所依据的行业景气、薪酬数据大多是第三方调研，读者需要自行甄别。
\end{enumerate}

\subsection{边界说明}
本调查不声称能回答“读完这个专业到底值不值”。这个问题最终取决于个人能力、学校层次、是否继续深造和市场行情的个人权衡。\textbf{本调查唯一能确认的是：供给侧确实有这个专业、规模不小、金融行业学历门槛有官方旁证，以及“量化人才供需的任何量化断言都没有官方依据”这一点。}

\section{本调查的局限与来源}

\subsection{局限}
\begin{itemize}[itemsep=2pt]
\item AMAC 的官方从业数据以 2022 年为最新可查（更新数据需上 AMAC 官网“统计数据”栏目）；这份调查不含 2023--2026 年的官方从业总数。
\item 教育部数据到 2025-12-31；阳光高考部分页面存在版本区间差异。
\item \textbf{最大的局限}：官方没有“量化”这个细分统计，所以本调查做不出任何“量化人才精确供需比”的结论——这是数据供给本身的客观限制，我们如实说明。
\end{itemize}

\subsection{读者如何自己核对}
想复核这份调查，可以访问以下一手来源：
\begin{itemize}
\item AMAC 官网统计数据 / 年报：\url{https://www.amac.org.cn/}
\item 教育部《本科专业目录（2025 年）》与年度备案审批公告：\url{http://www.moe.gov.cn/srcsite/a08/moe_1034/s4930/}
\item 阳光高考专业库：\url{https://gaokao.chsi.com.cn/zyk/zybk/}
\item 中国证券业协会研究（经证券时报）：\url{https://www.stcn.com/article/detail/3149883.html}
\item \small 中基协从业人员公示平台（逐家核算头部从业数）：\url{https://gs.amac.org.cn/amac-infodisc/res/pof/fund/index.html}
\end{itemize}

\section{关于非权威数据：读者常看到的数字，能信到什么程度}

本调查的正文只采用有官方或可靠出处、且口径可核验的数据。但读者在行业报道、招聘帖、自媒体里，还会频繁看到另一批数字——它们\textbf{人气很高，但出处不是一个官方发布的汇总}。这份调查把凡是出现过的、可能被误采信的非权威数据单独集中说明，方便读者辨识其性质和局限。

\begin{table}[ht]
\centering
\small
\begin{tabular}{@{}>{\raggedright\arraybackslash}p{3.4cm}>{\raggedright\arraybackslash}p{2.0cm}>{\raggedright\arraybackslash}p{2.4cm}>{\raggedright\arraybackslash}p{3.8cm}@{}}
\toprule
\textbf{数字} & \textbf{性质} & \textbf{出处} & \textbf{局限 / 说明} \\
\midrule
百亿量化私募从业约 4,495 人（74 家，2026-07） & 第三方统计 & 私募排排网 & 非官方；口径多版本并存（71/74/85 家）；本调查\textbf{额外采用}但明确标注 \\
机构量化团队约 1--2 万人 & 经验估算 & 行业口头 & 无统计出处，本调查\textbf{不采信} \\
个人全职量化交易约 10--20 万人 & 经验估算 & 行业口头 & 无统计出处，本调查\textbf{不采信} \\
合规缴纳社保量化从业约 3,000--4,500 人 & 推算值 & 本调查早年推算 & \textbf{已撤回}，无 AMAC 依据 \\
量化私募从业总数约 2.58 万人 & 第三方估算（2021） & 《2021 中国量化投资白皮书》 & 基于“量化占比”估算，时间早，口径假设强 \\
金融学类读研比例约 20\% & 第三方调研 & 麦可思《就业蓝皮书》 & 非教育部官方；按专业类无官方对应值 \\
量化从业者博士 11\%、硕士 65\% & 第三方调研 & 第三方白皮书引证 & 非官方，多为商业机构抽样 \\
\bottomrule
\end{tabular}
\caption{读者常见的非权威数据清单及说明}
\end{table}

\noindent\textbf{如何看待这张表}：
\begin{enumerate}[itemsep=2pt]
\item \textbf{“第三方”不一定错，但要知道它的层级}：表里的数字有些来自私募排排网、麦可思等有一定口碑的商业机构，可以作为\textbf{参考方向}；但它们的口径、抽样、时点各不相同，\textbf{不能当作官方事实引用}。
\item \textbf{“经验估算”和“推算值”应视为不可核验}：机构 1--2 万、个人 10--20 万、合规社保 3,000--4,500 等，要么是行业口头数字，要么建立在很强的假设上，本调查\textbf{不予采信}，列出仅为防读者误用。
\item \textbf{引用这些数字时的建议}：请务必注明“来源为某机构统计/估算、时点与口径”，并区分“这是某机构的口径”与“官方没有这个数”。
\end{enumerate}

主报告的结论（下一节）仍然只建立在官方可查证的数据之上；这张表是\textbf{独立的辨析附录}，不改变正文的权威性口径。

\section{结论}

\begin{enumerate}[itemsep=2pt]
\item \textbf{官方没有“量化从业人员”的细分统计}：任何“量化从业者 X 万人”都是第三方推算，不是权威数据。（这类常见的第三方数字及各自的权威层级，已在上一节《关于非权威数据》集中列出。）
\item \textbf{供给侧专业客观存在、规模不小}：金融工程年毕业约 2.2--2.4 万（教育部官方口径），是大规模培养方向。
\item \textbf{学历门槛有官方旁证}：券商从业者硕士占比 34.06\%（中证协）且在上升，提示“本科不深造直接就业”难度较高。
\item \textbf{“量化供需是否匹配”无法用官方数据回答}：一方面，官方两边都没有“量化”这个标签；另一方面，需求侧可得的是在岗存量（18.30 万），供给侧可得的是当年毕业生流量（约 2.5--2.8 万），\textbf{存量与流量口径不可比}，而又缺少官方的年度新增岗位数据作基准，因此任何量化供需比的结论都无法被证实。
\item \textbf{对“值不值得读”}：本调查给出分人群的判断（规模培养 + 学历内卷 + 深造更优），但明确说明它不对未来的薪资、就业做量化结论，第三方行业数据只能视为参考。
\item \textbf{对非权威数据的立场}：凡第三方统计、经验估算或推算值，本调查一律单独说明其出处与局限（见上节清单），其中多数不作结论依据；只有“头部从业约 4,495 人”因读者高度关注且无官方替代，才破例采用并明确标注。读者在别处看到类似数字时，应先问“这是谁的统计、口径与时点是什么”。
\end{enumerate}

\noindent\textbf{声明}：本调查坚持以官方和权威渠道的数据为准，第三方数据一律标注清楚。文中的判断仅反映在可查证数据范围内的有限结论，不构成升学或职业建议。

\end{document}